\documentclass[conference]{IEEEtran}
\IEEEoverridecommandlockouts

% ------------------------- Typography ---------------------------------
\usepackage[T1]{fontenc}
\usepackage[utf8]{inputenc}
\usepackage{newtxtext}      % Times-like text font
\usepackage{newtxmath}      % matching math font
\usepackage{microtype}      % subtle kerning/spacing polish

% ------------------------- Packages -----------------------------------
\usepackage{cite}
\let\Bbbk\relax
\usepackage{amsmath,amssymb,amsfonts}
\usepackage{algorithmic}
\usepackage{graphicx}
\usepackage{textcomp}
\usepackage{xcolor}
\usepackage{booktabs}       % professional horizontal rules for tables
\usepackage{array}
\usepackage{tabularx}
\usepackage{url}
\usepackage[hidelinks]{hyperref}

\newcommand{\comparisonTableSetup}{%
  \small
  \setlength{\tabcolsep}{3pt}%
  \renewcommand{\arraystretch}{1.25}%
}

\begin{document}

%======================================================================
% TITLE
%======================================================================
\title{A Comprehensive Survey on AI-Assisted Speech Therapy for Children with Autism Spectrum Disorder}

%======================================================================
% AUTHORS
%======================================================================
\newcommand{\authcol}[1]{%
  \begin{minipage}[t]{0.31\textwidth}\centering\small #1\end{minipage}}

\author{%
% ---------- Row 1 ----------
\authcol{%
  \textbf{Rohith Kumar Dhamagatla}\\[2pt]
  \textit{Department of Computer Science and Engineering}\\
  \textit{Amrita School of Computing, Coimbatore}\\
  \textit{Amrita Vishwa Vidyapeetham, India}\\
  cb.sc.u4cse23018@cb.students.amrita.edu}
\hfill
\authcol{%
  \textbf{Thungamitta Venkataramana}\\[2pt]
  \textit{Department of Computer Science and Engineering}\\
  \textit{Amrita School of Computing, Coimbatore}\\
  \textit{Amrita Vishwa Vidyapeetham, India}\\
  cb.sc.u4cse23055@cb.students.amrita.edu}
\hfill
\authcol{%
  \textbf{Ajay Bhargava Jashwanth Reddy Vustelamuri}\\[2pt]
  \textit{Department of Computer Science and Engineering}\\
  \textit{Amrita School of Computing, Coimbatore}\\
  \textit{Amrita Vishwa Vidyapeetham, India}\\
  cb.sc.u4cse23058@cb.students.amrita.edu}
\\[12pt]

\authcol{%
  \textbf{Challa Kalyan Kumar Reddy}\\[2pt]
  \textit{Department of Computer Science and Engineering}\\
  \textit{Amrita School of Computing, Coimbatore}\\
  \textit{Amrita Vishwa Vidyapeetham, India}\\
  cb.sc.u4cse20360@cb.students.amrita.edu}
\hfill
\authcol{%
  \textbf{Sholingram Hemanth}\\[2pt]
  \textit{Department of Computer Science and Engineering}\\
  \textit{Amrita School of Computing, Coimbatore}\\
  \textit{Amrita Vishwa Vidyapeetham, India}\\
  cb.sc.u4cse23504@cb.students.amrita.edu}
\hfill
\authcol{%
  \textbf{Dr. Venkataraman D}\\[2pt]
  \textit{Department of Computer Science and Engineering}\\
  \textit{Amrita School of Computing, Coimbatore}\\
  \textit{Amrita Vishwa Vidyapeetham, India}\\
  d\_venkataraman@cb.amrita.edu}
}

\maketitle

%======================================================================
% ABSTRACT
%======================================================================
\begin{abstract}
Speech and language difficulties continue to be a concern for caregivers and healthcare professionals. This is due to the fact that these issues frequently require monitoring and addressing over a period. In recent years, artificial intelligence has significantly influenced speech therapy research. It aids in speech evaluation by identifying communication issues promptly and facilitating intelligent rehabilitation methods. This study examines various artificial intelligence methods, such as machine learning, deep learning, speech recognition, and additional techniques. These methods assist in identifying speech issues, recognizing children's speech, analyzing emotional expression in speech, and supporting individuals with Autism Spectrum Disorder. By concentrating on a single technique, this research categorizes the current literature into key areas, compares the approaches employed by various studies, and discusses their advantages and disadvantages. The comprehensive evaluation indicates that intelligent speech therapy systems have improved speech assessment efficiency and enabled greater access to rehabilitation for many individuals. Nonetheless, issues such as restricted language resources, varying evaluation approaches, privacy issues, and computational complexity remain unresolved. The results indicate that future advancements rely on employing intelligence methods that are user-friendly, energy-efficient, and capable of handling various types of data.
\end{abstract}
%======================================================================
% I. INTRODUCTION
%======================================================================
\section{Introduction}

Human interaction relies heavily on spoken language, making speech disorders a critical concern. Conditions such as Autism Spectrum Disorder (ASD), dysarthria, apraxia, and developmental speech disorders frequently require sustained intervention. While conventional speech therapy benefits many individuals, access remains limited and costly, motivating Artificial Intelligence (AI)-driven systems that offer scalable, patient-centered alternatives. Research remains dispersed across domains, making comprehensive understanding difficult. This survey addresses this gap by examining AI applications in speech therapy. The key contributions are:

\begin{itemize}
    \item Reviewing AI advances in speech therapy across six thematic categories.
    \item Comparing methodologies, datasets, and evaluation approaches across studies.
    \item Identifying research obstacles including data scarcity, multilingual support, and explainability.
    \item Highlighting future opportunities for speech therapy system development.
\end{itemize}

Section~II presents background concepts. Section~III reviews and classifies current studies. Section~IV provides comparative analysis. Section~V discusses open issues. Section~VI concludes the survey.

%======================================================================
% II. BACKGROUND / PRELIMINARIES
%======================================================================
\section{Background / Preliminaries}
\label{sec:background}

\subsection{Autism Spectrum Disorder and Speech}

ASD is a neurodevelopmental condition involving social communication deficits and repetitive behaviors. Speech abnormalities---atypical prosody, rhythm, and vocalization---emerge early and parallel typical speech development, motivating research into speech as a diagnostic and severity assessment tool~\cite{ref7,ref10}.

\subsection{Acoustic and Prosodic Indicators}

Research has identified acoustic signatures of ASD, including prosodic features in preschoolers~\cite{ref7}, automated Vowel Space Area (VSA) computation for articulatory accuracy~\cite{ref8}, and correlations between speech and fine motor skills linked through shared motor control systems~\cite{ref6}.

\subsection{Speech Disorders Beyond ASD}

Related speech motor disorders including dysarthria exhibit similar deterioration in intelligibility and articulation. These disorders have been examined using wav2vec embeddings~\cite{ref4,ref20}, which transfer to ASD speech analysis due to analogous signal processing challenges.

\subsection{Deep Learning and Self-Supervised Representations}

Self-Supervised Learning (SSL), particularly wav2vec2, underpins much current work. Wav2vec2 has been adapted for child speech recognition under limited labeled data~\cite{ref1}. SSL embeddings have been combined with extended Geneva Minimalistic Acoustic Parameter Set (eGeMAPS) features for ASD detection~\cite{ref2}. Transformer transfer learning has been applied for ASD recognition~\cite{ref3}, while cross-lingual and phoneme-based transfer has been enabled for low-resource settings~\cite{ref5,ref22}.

\subsection{Assistive and Therapeutic Technologies}

Beyond diagnosis, technologies include Augmentative and Alternative Communication (AAC) devices~\cite{ref17}, digital games~\cite{ref12}, mobile apps~\cite{ref13}, deep learning-based mobile therapy~\cite{ref14}, generative AI platforms~\cite{ref15}, and Opti-Speech's 3D visual feedback~\cite{ref16}. Screening approaches combine Automatic Speech Recognition (ASR)~\cite{ref9} with multimodal AI~\cite{ref10}.

\subsection{Privacy and Edge Deployment}

Given the sensitivity of children's speech data, federated learning has been applied to keep audio on-device~\cite{ref24}, while Tiny Machine Learning (TinyML) speech recognition via model quantization has been demonstrated for edge deployment~\cite{ref23}.

%======================================================================
% III. LITERATURE SURVEY / TAXONOMY
%======================================================================
\section{Literature Survey / Taxonomy}
\label{sec:survey}

We organize the surveyed papers into six thematic categories reflecting the layers of an AI-enabled speech therapy solution: (1)~representation learning, (2)~biomarkers, (3)~multimodal screening and intervention, (4)~interactive delivery, (5)~affective computing, and (6)~enabling infrastructure. Table~\ref{tab:comparison} provides a comparative summary. A unifying constraint---labeled autistic child speech---affects all layers, driving the field toward transfer learning, data-efficient methods, and privacy-preserving techniques. The themes form a pipeline: limitations at one layer define the agenda of the next.

%----------------------------------------------------------------------
\subsection{Self-Supervised Speech Representations for Child and Atypical Speech}

This theme collects studies adapting large SSL models, primarily wav2vec~2.0, to child and atypical speech where labeled data is scarce. All share the assumption that representations learned from abundant unlabeled data can transfer to target applications with minimal labeled data.

Extensive ablation studies over pretraining and fine-tuning combinations for child ASR found that ten hours of child speech during fine-tuning sufficed to outperform baselines (Word Error Rate (WER): 7.42 on MyST, 2.91 on PF-STAR, 12.77 on CMU~Kids), though recognition of the youngest children remained problematic~\cite{ref1}. ASD detection from infant voice was targeted by combining wav2vec2.0 contextual vectors with an eGeMAPS-based autoencoder in a Bidirectional Long Short-Term Memory (BLSTM) classifier, improving accuracy from 64.74\% to 71.66\% and Unweighted Average Recall (UAR) from 65.04\% to 70.81\%~\cite{ref2}.

The Noor Project compared Discriminative Fine-Tuning (D-FT) with wav2vec2.0 Fine-Tuning, achieving 94.8\% UAR on binary classification versus 91.5\%, but performance dropped sharply on four-way diagnosis (60.9\% vs.\ 54.3\%), quantifying how limited labels impair fine-grained diagnosis~\cite{ref3}. Wav2vec applied to dysarthria revealed that first-layer embeddings best served detection (+1.23\% accuracy) while final-layer embeddings excelled at severity classification (+10.62\%)~\cite{ref4}. Cross-lingual transfer was addressed through knowledge distillation, reducing the transfer gap for low-resource languages on CoVoST-2~\cite{ref5}.

\emph{Synthesis.} SSL pretraining with wav2vec2.0 is established as the preferred approach for data-sparse child speech~\cite{ref1,ref2,ref3,ref4,ref5}. Common findings include the effectiveness of data-efficient fine-tuning, task-appropriate layer selection, and the persistent adult-to-child domain gap.

A particularly instructive technical pattern emerges from comparing layer-wise findings~\cite{ref4} with dual-stream architectures~\cite{ref2}. Shallow transformer layers capture phonetic surface features suited for binary detection, while deeper layers encode severity-level abstractions---a gradient that maps directly onto the clinical workflow, where a therapist first identifies the presence of an error and then grades its severity. A complementary conclusion shows that pairing an interpretable eGeMAPS stream with an opaque SSL stream outperforms either alone, suggesting that the optimal architecture for child speech applications is not a single end-to-end model but a hybrid that lets clinicians inspect one channel while the other compensates for the interpretable channel's ceiling. Together, these observations imply that future therapy systems should expose multiple representation depths simultaneously, enabling both coarse screening and fine-grained pronunciation feedback within a single inference pass.

%----------------------------------------------------------------------
\subsection{Acoustic and Prosodic Biomarkers of ASD}

This theme examines interpretable acoustic features revealing ASD, complementing SSL embeddings with clinically meaningful markers.

Acoustic, video, and handwriting data from five ASD and five control children were correlated, finding distinct motor coordination patterns supporting the theory linking ASD speech difficulties to motor coordination problems~\cite{ref6}. This approach was scaled to 74 preschoolers with a custom diarization pipeline, identifying prosodic biomarkers predictive of developmental outcomes up to a year ahead~\cite{ref7}. VSA estimation was automated, transforming a labor-intensive measurement into a scalable articulation metric~\cite{ref8}. Convolutional Neural Network (CNN)-derived features were used for analyzing autistic children's voices, leveraging learned representations to handle high vocal expression variance~\cite{ref9}.

\emph{Synthesis.} Measurable ASD fingerprints across motor coordination, prosody, vowels, and speaker acoustics have been demonstrated~\cite{ref6,ref7,ref8,ref9}, providing interpretable criteria that can audit black-box predictions. Small, heterogeneous datasets remain a common limitation.

From a clinical translation perspective, the prognostic finding that certain prosodic patterns in non-verbal preschoolers predicted communicative outcomes a full year later carries significant implications for therapy system design~\cite{ref7}. If a computational biomarker can forecast developmental trajectories, a therapy platform can move beyond reactive correction toward proactive goal-setting, adjusting intervention intensity before a child falls behind rather than after. Similarly, the automated VSA metric offers a repeatable, objective outcome measure that could replace subjective clinician ratings in longitudinal studies, addressing the evaluation inconsistency highlighted in Section~\ref{sec:future}~\cite{ref8}. The motor-coordination framework, though limited by sample size, opens a pathway to multisensory assessment: if speech and handwriting share a motor substrate in ASD, then a tablet-based therapy app could simultaneously collect articulatory audio and stylus kinematics, doubling the diagnostic signal without additional hardware~\cite{ref6}.

%----------------------------------------------------------------------
\subsection{Multimodal AI for Screening and Intervention}

This theme marks the transition from passive measurement to active decision-making systems that combine multiple data streams for clinical action.

A population-wide two-stage multimodal system was developed using data from 1,242 children aged 18--48 months captured via a mobile application~\cite{ref10}. Rather than relying solely on questionnaire scores, the system performs Natural Language Processing (NLP) on the text of screening questions (Modified Checklist for Autism in Toddlers (M-CHAT), Social Communication Questionnaire-Lifetime (SCQ-L), Social Responsiveness Scale (SRS)) to extract behavioural markers, combined with audio features from parent--child interactions. Stage~1 classified children into typically developing and high-risk/ASD categories with Area Under the Receiver Operating Characteristic Curve (AUROC) 0.942; Stage~2 discriminated high-risk from ASD with AUROC 0.914 and 85.2\% accuracy. Predicted risk categories showed 79.59\% concordance with gold-standard Autism Diagnostic Observation Schedule, Second Edition (ADOS-2) evaluations and high correlation ($r=0.830$, $p<0.001$), providing strong evidence that mobile-collected multimodal data can approximate expert-level screening.

The loop from diagnosis to personalized intervention was closed in a subsequent study~\cite{ref11}. A Deep Neural Network (DNN) achieved 96.98\% accuracy, 97.65\% precision, 96.74\% recall, and 99.75\% AUROC, outperforming Random Forest and Logistic Regression baselines. Feature selection via Least Absolute Shrinkage and Selection Operator (LASSO) and Random Forest identified key predictors including Quantitative Checklist for Autism in Toddlers (Q-CHAT-10) scores. A Deep Deterministic Policy Gradient (DDPG) reinforcement-learning agent then simulated intervention policies over 12 months, varying therapy types, frequencies, and intensities. Simulations showed up to 25\% improvement in social abilities, 30\% reduction in behavioural problems, and high-risk subjects dropping from 65\% to 25\%, though these remain virtual outcomes requiring longitudinal clinical validation.

\emph{Synthesis.} The progression from one-shot categorization to scalable screening and closed-loop personalized intervention is captured by these studies~\cite{ref10,ref11}. The former leverages mobile sensor-based data fusion, the latter reinforcement learning-based therapy planning; together they define the adaptivity needed for a comprehensive therapy platform. The shared caveat---reliance on single-site data and simulated outcomes---frames the validation challenges discussed in Section~\ref{sec:future}.

%----------------------------------------------------------------------
\subsection{Digital Games, Mobile Apps, and Interactive Platforms}

This theme addresses the application layer where therapy meets the child, encompassing systematic reviews of consumer technology and systems papers constructing interactive platforms.

Two systematic reviews set the context. An analysis of 27 computer games for child speech therapy across four databases found Mel-Frequency Cepstral Coefficients (MFCC) features with PocketSphinx as the most common recognition pipeline and 48\% of games targeting parent-assisted practice~\cite{ref12}. While gamification raised satisfaction and motivation, key barriers included frustration from ASR errors, noisy environments, and inappropriate difficulty levels. An evaluation of over 5,000 mobile apps using the Mobile Application Rating Scale (MARS) found only 3\% met criteria for analysis and none rated excellent---exposing a significant gap between application quantity and therapeutic quality~\cite{ref13}.

System papers show advancing delivery capabilities. A TensorFlow/ARKit mobile platform was created integrating image classification, text-to-speech, NLP, and cloud storage for home-based ASD rehabilitation with clinician data sharing~\cite{ref14}. A generative AI tablet platform was tested in a pre/post study with 22 children aged 2--12, collecting 2,000 audio recordings and finding significant language improvement for children engaged over 15 hours~\cite{ref15}. Opti-Speech, a real-time 3D visual feedback system using Electromagnetic Articulography (EMA) to monitor tongue and jaw movements within a transparent head model, provides high-fidelity articulation training limited to clinical settings by specialized hardware~\cite{ref16}. A review of AAC aids using Tau-U and improvement-rate difference analysis found high-tech devices (Tau-U 0.80--1.00) significantly more effective than low-tech alternatives for social communication and speech production~\cite{ref17}.

\emph{Synthesis.} The engagement and delivery layer is provided by these studies~\cite{ref12,ref13,ref14,ref15,ref16,ref17}, and AAC evidence confirms that technology-rich interfaces yield real communication benefits~\cite{ref17}. However, both systematic reviews reach the same conclusion~\cite{ref12,ref13}: creative interfaces are plentiful, yet reliable child speech recognition and rigorous efficacy evidence remain insufficient---the very deficiency that the SSL, biomarker, and enabling technology themes must jointly overcome.

An important cross-theme observation connecting the delivery and affective layers is the question of \emph{when} and \emph{how} feedback should be delivered to maintain a child's engagement without inducing frustration. The catalogued barriers---lowered self-esteem from repeated failure, inappropriate difficulty levels, and ASR-induced false negatives---are essentially affective problems masquerading as technical ones~\cite{ref12}. A child who receives incorrect negative feedback from a faulty recognizer does not simply lose a game turn; the child experiences a violation of trust in the system, which erodes future willingness to engage. Conversely, children gravitate toward high-tech AAC devices precisely because richer feedback sustains motivation~\cite{ref17}. This suggests that the affect-recognition capabilities reviewed in Theme~5 are not a separate luxury but a necessary control loop for the delivery mechanisms of Theme~4: an affect-aware system could suppress uncertain ASR judgments when it detects rising frustration, or increase challenge when it senses sustained engagement, thereby closing the gap between interface creativity and therapeutic reliability that both systematic reviews identified.

%----------------------------------------------------------------------
\subsection{Speech Emotion and Affect Recognition}

Affect recognition is indispensable for maintaining autistic children's attention: a therapy system must detect frustration, boredom, and engagement to adapt its responses, and emotional expression is itself atypical in ASD, making affect both a control signal and a treatment target.

A comprehensive review of multimodal Speech Emotion Recognition (SER) fusion approaches---early, late, deep, and hybrid---alongside attention models and graph-based fusion, demonstrated that unimodal voice-only methods break under noise and ambiguity, precisely the conditions of home therapy~\cite{ref18}. The resulting fusion typology provides a framework for combining voice, video, and interaction data within a child's therapy session, compatible with the multimodal screening architecture of Theme~3. A Preferred Reporting Items for Systematic Reviews and Meta-Analyses (PRISMA)-informed systematic review of automatic emotion recognition in autistic children examined facial expressions, speech prosody, and physiological measures~\cite{ref19}. Support Vector Machines (SVMs) and neural networks were found to be leading classifiers, basic emotions (happiness, fear, sadness) most recognized, single-channel architectures preferred, and early fusion dominating over late fusion, while highlighting both the lack of annotated emotion datasets for autistic children and the practical challenge of modality disappearance (e.g., face turning away) in real-life settings.

\emph{Synthesis.} Affect recognition can provide real-time feedback signals for adaptive therapy systems~\cite{ref18,ref19}. The fusion approaches complement the multimodal screening of Theme~3, suggesting a unified fusion framework for both screening and engagement. The recurring data scarcity for children's emotional expression mirrors findings across all themes.

%----------------------------------------------------------------------
\subsection{Enabling Technologies: Enhancement, Alignment, and Edge/Privacy Deployment}

This final theme assembles the infrastructure needed to make the preceding systems reliable and practical, spanning field synthesis, signal enhancement, phonetic alignment, on-device inference, and privacy-preserving training.

The broadest scope was provided by a synthesis of 231 papers (2007--2024) on Machine Learning (ML)/Deep Learning (DL) for sixteen speech disorders, tracking progression from SVMs to hybrid neural architectures and discussing the evolution from acoustic to prosodic and linguistic features~\cite{ref20}. The study identifies fundamental field-wide limitations---lack of longitudinal and cross-cultural data, and insufficient clinical interpretability of AI models---and recommends Large Language Models (LLMs) and adaptive approaches as paths toward disorder-agnostic, personalized medicine. Six neural vocoders were compared against two phase-reconstruction baselines using Multiple Stimuli with Hidden Reference and Anchor (MUSHRA) listening tests, evaluating quality against computational cost to provide practitioners concrete guidelines for selecting enhancement modules where therapy systems must restore or generate child speech under adverse conditions~\cite{ref21}.

Contrastive Language-Audio Pretraining with International Phonetic Alphabet (CLAP-IPA), a contrastive embedding model trained on IPAPACK (115 languages), was developed to enable open-vocabulary phoneme-speech matching and zero-shot forced alignment across 95 unseen languages without retraining---the technical foundation for syllable-level pronunciation feedback across diverse linguistic settings~\cite{ref22}. A quantized one-dimensional CNN (1D-CNN) was trained on an Arduino Nano~33 BLE Sense microcontroller, achieving 97\% accuracy for 23 keywords and demonstrating that real-time, offline, privacy-preserving inference is feasible on TinyML hardware~\cite{ref23}. Federated learning with secure aggregation and Generative Adversarial Network (GAN)-generated imposter samples was applied, keeping audio data on-device while maintaining speaker recognition accuracy competitive with centralized training~\cite{ref24}.

\emph{Synthesis.} All pipeline components are connected by this theme~\cite{ref20,ref21,ref22,ref23,ref24}: from the field map and explainability framework, through enhancement and synthesis, to the alignment system enabling syllable-level feedback, and the edge and federated deployment systems ensuring private, offline, and personalized operation.

%----------------------------------------------------------------------
\subsection{Cross-Category Synthesis}

Viewed collectively, the six themes trace a single vertical stack in which each layer depends on those beneath and is constrained by them. Dependable child-facing delivery~\cite{ref12,ref13,ref14,ref15,ref16,ref17} presupposes reliable perception and scoring, which in turn require robust representations~\cite{ref1,ref2,ref3,ref4,ref5} and interpretable biomarkers~\cite{ref6,ref7,ref8,ref9}. Adaptive, affect-aware personalization~\cite{ref10,ref11,ref18,ref19} assumes both of those layers, and every layer assumes the enabling infrastructure~\cite{ref20,ref21,ref22,ref23,ref24} that keeps it deployable and private. A single constraint---annotated speech from autistic children---throttles all six layers simultaneously, making the privacy-preserving and data-efficient approaches of Themes~1 and~6 not optional add-ons but the very mechanisms by which the field can break that bottleneck. This taxonomy therefore does more than classify papers: it reveals how gains at one layer propagate through the entire pipeline, and it is this structure upon which our comparative analysis and research agenda rest.

%======================================================================
% IV. COMPARATIVE ANALYSIS & DISCUSSION
%======================================================================
\section{Comparative Analysis and Discussion}
\label{sec:analysis}

This section presents a critical comparison of the surveyed studies. Several methodological trends emerge from the collective analysis. First, there is a clear progression from traditional classifiers (SVMs, Random Forests) toward self-supervised transformer architectures, with wav2vec2.0 serving as the dominant backbone across recognition, detection, and classification tasks. Second, modality fusion is expanding from single-stream audio processing toward multimodal pipelines that integrate voice, text, video, and physiological signals---a shift driven by the recognition that unimodal approaches are insufficient for the noisy, variable conditions of real-world child interaction. Third, the field is moving beyond static classification toward closed-loop intervention systems, exemplified by reinforcement learning-based therapy planning~\cite{ref11}. Finally, a persistent tension exists between model sophistication and practical deployment: the most accurate systems often demand computational resources incompatible with the edge devices and privacy constraints that clinical applications require. Tables~\ref{tab:comparative_analysis_part1} and~\ref{tab:comparative_analysis_part2} summarize each study's methodology, datasets, metrics, key findings, and limitations.

\clearpage
\onecolumn

\begin{table}[p]
\centering
\caption{Comparative Analysis of Surveyed Studies (Part~1)}
\label{tab:comparative_analysis_part1}
\label{tab:comparison}
\comparisonTableSetup
\begin{tabularx}{\textwidth}{|>{\hsize=0.20\hsize\raggedright\arraybackslash}X|
                              >{\hsize=1.20\hsize\raggedright\arraybackslash}X|
                              >{\hsize=1.10\hsize\raggedright\arraybackslash}X|
                              >{\hsize=0.90\hsize\raggedright\arraybackslash}X|
                              >{\hsize=0.70\hsize\raggedright\arraybackslash}X|
                              >{\hsize=1.40\hsize\raggedright\arraybackslash}X|
                              >{\hsize=1.50\hsize\raggedright\arraybackslash}X|}
\hline
\textbf{No.} & \textbf{Paper Title} & \textbf{Method} & \textbf{Dataset} & \textbf{Metrics} & \textbf{Key Findings} & \textbf{Limitations} \\
\hline
1 & Cross-lingual Transfer Learning for Low-Resource Speech Translation (2024) & SAMU-XLS-R + Knowledge Distillation & CoVoST-2 (128 languages) & BLEU-4 & Closes cross-lingual transfer gap; improves low-resource results significantly. & Full encoder fine-tuning causes catastrophic forgetting; adapters required. \\
\hline
2 & Clinical Efficacy of AAC in Minimally Verbal Children with ASD (2023) & Systematic Review (AAC) & 8 Clinical Studies (RCTs, SCEDs) & Tau-U, IRD, PEDro & High-tech AAC surpasses low-tech for social communication in ASD. & Few RCTs; small sample sizes available. \\
\hline
3 & Multimodal Fusion in Speech Emotion Recognition (2026) & Systematic Review (Multimodal SER) & IEMOCAP, MELD, RAVDESS & WA, UA & Hybrid and graph-based fusions optimal for cross-modal emotion detection. & Heavy computational burden; modality misalignment errors. \\
\hline
4 & Digital Games for Speech Therapy in Children (2022) & Systematic Review (Digital Games) & 27 Articles (PubMed, Scopus) & User Feedback, ASR Accuracy & Gamification increases engagement in articulation therapy. & ASR errors cause frustration; highly noise-sensitive. \\
\hline
5 & Wav2vec2 Self-Supervised Learning for Child Speech Recognition (2023) & Wav2vec2 SSL Pre-training & MyST, PFSTAR, CMU\_KIDS & WER & Fine-tuning on small child speech samples yields SOTA recognition performance. & Adult--child domain mismatch; large models offer no benefit over base. \\
\hline
6 & AI for Classification and Enhancement of Speech/Language Disorders (2024) & Systematic Review (AI for Speech Disorders) & 231 Studies (TORGO, UA-Speech) & Accuracy, WER, STOI, PESQ & Deep learning enhances disorder diagnosis and signal quality. & Lacks longitudinal and cross-cultural data; low AI clinical interpretability. \\
\hline
7 & Automated Assessment of Vowel Space Area (2013) & Automated VSA via convex hull & TIMIT Corpus & Pearson Correlation & Automated VSA precisely measures total articulatory workspace. & Requires phonetically balanced speech for valid comparison. \\
\hline
8 & Speech and Fine Motor Coordination in ASD Children (2020) & Cross-correlations + multivariate GMMs & 10 Children (5 ASD, 5 NT) & AUC, Cohen's~d & Speech--motor cross-correlations distinguish ASD from NT controls. & Tiny pilot sample; needs larger validation groups. \\
\hline
9 & Treatment Apps for Speech Disorders in Children (2018) & Systematic Appraisal (MARS) & 132 Speech Therapy Apps & MARS Scores, ICC & Apps have decent functionality but lack therapeutic validity. & English-only; no real therapeutic efficiency correlation. \\
\hline
10 & Federated Learning for Speaker Recognition (2021) & Federated learning + secure aggregation & VoxCeleb-1 & EER, t-test & Preserves privacy without sacrificing recognition accuracy. & High bandwidth needed; secure aggregation slightly reduces performance. \\
\hline
11 & ASR for ASD Analysis in Young Children (2022) & CNN, Random Forest, MFCC, NLP & Autism Kids' Speech Corpus & UAR, Accuracy & CNN features match traditional approaches for ASD speech classification. & Computationally expensive; dynamic thresholding may be inaccurate. \\
\hline
12 & Comparison of Neural Vocoders for Speech Reconstruction (2019) & Neural Vocoders vs.\ Phase Reconstruction & LJ Speech Corpus & MUSHRA listening test & Autoregressive models provide outstanding perceptual quality. & Phase reconstruction yields low scores due to pseudo-inverse issues. \\
\hline
\end{tabularx}
\end{table}

\clearpage

\begin{table}[p]
\centering
\caption{Comparative Analysis of Surveyed Studies (Part~2)}
\label{tab:comparative_analysis_part2}
\comparisonTableSetup
\begin{tabularx}{\textwidth}{|>{\hsize=0.20\hsize\raggedright\arraybackslash}X|
                              >{\hsize=1.20\hsize\raggedright\arraybackslash}X|
                              >{\hsize=1.10\hsize\raggedright\arraybackslash}X|
                              >{\hsize=0.90\hsize\raggedright\arraybackslash}X|
                              >{\hsize=0.70\hsize\raggedright\arraybackslash}X|
                              >{\hsize=1.40\hsize\raggedright\arraybackslash}X|
                              >{\hsize=1.50\hsize\raggedright\arraybackslash}X|}
\hline
\textbf{No.} & \textbf{Paper Title} & \textbf{Method} & \textbf{Dataset} & \textbf{Metrics} & \textbf{Key Findings} & \textbf{Limitations} \\
\hline
13 & End-to-End ASD Detection Using Pre-trained Model (2022) & Wav2vec2.0 + BLSTM & SNUBH data, Living Lab (191 children) & Accuracy, F1, UAR & E2E joint optimization greatly improves ASD detection over autoencoder. & Requires manual speech segmentation; vulnerable to overlapping speech. \\
\hline
14 & Deep Mobile Linguistic Therapy for ASD (2022) & TensorFlow, ARKit, NLP & CALTECH-256, Aristo Mini & Training/Val.\ Acc./Loss & Mobile platform makes speech therapy more accessible. & English-only; severe disorders require additional assistance. \\
\hline
15 & Opti-Speech: Real-time 3D Visual Feedback (2014) & EMA tracking + 3D visual feedback & 4 Adult English Speakers & Target accuracy & Accurate real-time visual feedback on place of articulation. & No individual anatomy model; physical tongue sensors required. \\
\hline
16 & Multimodal AI for ASD Risk Stratification (2025) & Two-stage AI (RoBERTa, Whisper) & 1,242 children (18--48 mo.) & AUROC, Pearson corr. & Audio + screening texts effectively stratify early ASD risk. & Based on ADOS-2 not final clinical diagnosis; limited sample. \\
\hline
17 & The Noor Project: Fair Transfer Learning for ASD (2025) & D-FT \& Wav2Vec2 & De-Enigma, CPSD & UAR, F1-score & D-FT outperforms Wav2Vec2 in typicality and diagnostic classification. & Multi-class struggles with limited data; gender bias toward males. \\
\hline
18 & Automatic Emotion Recognition in Children with Autism (2022) & Systematic Review (PRISMA) & 50 qualitative, 27 quantitative studies & Modality/algorithm frequencies & SVMs and NNs most used; early fusion most common for ASD emotion detection. & Small, biased samples; ground truth labeling issues. \\
\hline
19 & GenAI Digital Platform for Language Therapy (2025) & Generative AI interactive therapy & 22 children (2--12 yrs) & Total words, TTR, MLU & Significant improvement for children engaged $>$15 hours of usage. & Small pilot; non-uniform groups; no control group. \\
\hline
20 & DNN and DDPG for Early ASD Diagnosis and Intervention (2025) & DNN + DDPG reinforcement learning & Kaggle ASD Datasets & Accuracy (96.98\%), ROC-AUC & DNN excels at screening; DDPG optimizes personalized intervention policies. & Dependent on virtual intervention data; needs longitudinal validation. \\
\hline
21 & TinyML for Speech Recognition (2025) & Quantized 1D CNN on IoT edge & Custom dataset (1 hr audio) & Accuracy, F1-Score & 97\% accuracy for 23 keywords on edge device. & Single-speaker dataset; auto-generated noise limits generalization. \\
\hline
22 & Open-Vocabulary Keyword Spotting via CLAP-IPA (2024) & CLAP-IPA + IPA-ALIGNER & IPAPACK (115 languages) & Hit@1, mAP, EER, F1 & Strong cross-lingual generalization for keyword spotting and alignment. & Compute-heavy for mobile; biased toward high-resource languages. \\
\hline
23 & Wav2vec-Based Dysarthria Detection and Classification (2023) & Wav2vec 2.0 + SVM classifiers & UA-Speech & Accuracy, F1-score & First-layer embeddings for detection; final-layer excels at severity grading. & Needs cross-database generalization and other disorder applicability study. \\
\hline
24 & Prosodic Signatures of ASD Severity in Preschoolers (2023) & PLSC + diarization + GeMAPS features & Geneva Cohort (74 preschoolers) & Bootstrap Ratios, p-values & Voice-quality correlates with ASD severity; prosodic patterns have prognostic value. & No TD control group; diarization optimized only for male examiners. \\
\hline
\end{tabularx}
\end{table}

\clearpage
\twocolumn

%======================================================================
% V. OPEN ISSUES & FUTURE RESEARCH DIRECTIONS
%======================================================================
\section{Open Issues and Future Research Directions}
\label{sec:future}

\subsection{Scarcity of Labeled Child Speech Data}

\textbf{Problem:} The limited availability of labeled speech data from autistic children remains a major challenge. Jain et al.~\cite{ref1} combined adult and child recordings, while the Noor Project~\cite{ref3} lacked sufficient samples for certain categories. Small and imbalanced datasets can reduce model accuracy, increase overfitting, and limit generalization across different children and speech conditions.

\textbf{Solution:} Future research should focus on collaborative data collection and the creation of larger datasets covering diverse age groups, languages, and speech conditions. Data augmentation, balanced sampling, transfer learning, and synthetic speech generation can help address data scarcity. Standardized datasets would also enable more reliable comparison between different AI-based approaches.

\subsection{Manual Preprocessing Bottleneck}

\textbf{Problem:} Many existing systems depend on time-consuming manual preprocessing and annotation. Lee et al.~\cite{ref2} manually segmented audio recordings, while Godel et al.~\cite{ref7} developed specialized diarization tools. Such manual processes require significant effort and are difficult to scale for large datasets and real-time applications.

\textbf{Solution:} Automated preprocessing pipelines integrating speech enhancement, voice activity detection, noise reduction, and speaker diarization should be developed. Self-supervised learning techniques could further reduce dependence on manually labeled data. These improvements would support scalable and real-time speech therapy applications.

\subsection{Insufficient Multimodal Integration}

\textbf{Problem:} Most studies primarily analyze speech, although autism assessment also involves facial expressions, eye contact, gestures, and body posture. Nguyen et al.~\cite{ref18} highlighted the limitations of voice-only methods in noisy environments, while Bae et al.~\cite{ref10} demonstrated the benefits of combining multiple data sources. Therefore, unimodal approaches may fail to capture important behavioral information.

\textbf{Solution:} Future systems should integrate speech, text, facial expressions, gestures, and physiological signals using multimodal learning architectures. Advanced fusion techniques could identify relationships between different behavioral and speech features. Such approaches may improve assessment accuracy and enable more personalized therapy.

\subsection{Personalized and Animated Feedback}

\textbf{Problem:} Personalized therapy and adaptive feedback mechanisms remain underexplored. Zheng et al.~\cite{ref11} applied reinforcement learning for intervention optimization, while Katz et al.~\cite{ref16} demonstrated the effectiveness of animated visual feedback. However, limited research has combined adaptive AI models with engaging feedback mechanisms for autistic children.

\textbf{Solution:} Future research should combine reinforcement learning with interactive animated feedback to dynamically personalize therapy exercises. Systems could adjust difficulty levels, feedback frequency, and learning activities according to each child's performance. This approach could improve engagement and support continuous therapy progress.

\subsection{Privacy and Edge Deployment}

\textbf{Problem:} Children's speech and behavioral data require strong privacy protection, while advanced AI models often demand significant computational resources. Existing privacy-preserving and edge-based approaches remain limited~\cite{ref23,ref24}. These limitations create challenges for deploying AI-based therapy systems on smartphones and other resource-constrained devices.

\textbf{Solution:} Lightweight models, federated learning, differential privacy, pruning, and quantization should be explored to enable secure and efficient processing. Performing AI computations directly on edge devices could reduce the transmission of sensitive information. Such techniques can improve privacy while supporting practical real-world deployment.

\subsection{Evaluation Gaps and Standards}

\textbf{Problem:} Existing studies lack standardized evaluation methods and sufficient real-world clinical validation. Differences in terminology, preprocessing, datasets, and performance metrics make reliable comparisons difficult~\cite{ref1,ref13,ref15,ref19}. Furthermore, many studies focus primarily on technical performance rather than measurable long-term clinical outcomes.

\textbf{Solution:} Future studies should establish common benchmark datasets, standardized performance metrics, and clinical evaluation frameworks. Evaluations should include diverse populations across different age groups, languages, and speech conditions. Collaboration between AI researchers and healthcare professionals could help ensure that developed systems provide measurable real-world therapy benefits.


%======================================================================
% VI. CONCLUSION
%======================================================================
\section{Conclusion}

This survey highlights AI's growing role in speech disorder assessment and therapy through machine learning, deep learning, and speech recognition. These techniques have enhanced speech analysis accuracy and enabled personalized, accessible rehabilitation. However, no single approach suits all contexts---dataset quality, language diversity, computational resources, and clinical validation critically influence performance. Future work must prioritize clinically validated, multimodal systems integrating privacy preservation, edge-appropriate models, and explainable AI. Sustained collaboration among researchers, clinicians, developers, and speech-language professionals is essential to transform these technologies into widely deployable solutions for diverse speech and communication disorders.

%======================================================================
% REFERENCES  (IEEE format, numbered in order of appearance)
%======================================================================
\begin{thebibliography}{00}

\bibitem{ref1} R. Jain, A. Barcovschi, M. Y. Yiwere, D. Bigioi, P. Corcoran, and H. Cucu, ``A wav2vec2-based experimental study on self-supervised learning methods to improve child speech recognition,'' \textit{IEEE Access}, vol.~11, pp.~46938--46948, 2023.

\bibitem{ref2} J. H. Lee, G. W. Lee, G. Bong, H. J. Yoo, and H. K. Kim, ``End-to-end model-based detection of infants with autism spectrum disorder using a pretrained model,'' \textit{Sensors}, vol.~23, no.~1, Art. no.~202, 2023.

\bibitem{ref3} N. D. Al Futaisi, B. W. Schuller, F. Ringeval, and M. Pantic, ``The Noor Project: Fair transformer transfer learning for autism spectrum disorder recognition from speech,'' \textit{Front. Digit. Health}, vol.~7, Art. no.~1274675, 2025.

\bibitem{ref4} F. Javanmardi, S. Tirronen, M. Kodali, S. R. Kadiri, and P. Alku, ``Wav2vec-based detection and severity level classification of dysarthria from speech,'' in \textit{Proc. IEEE Int. Conf. Acoust., Speech Signal Process. (ICASSP)}, Rhodes Island, Greece, 2023, pp.~1--5.

\bibitem{ref5} S. Khurana, N. Dawalatabad, A. Laurent, L. Vicente, P. Gimeno, V. Mingote, and J. Glass, ``Cross-lingual transfer learning for low-resource speech translation,'' arXiv:2306.00789, 2024.

\bibitem{ref6} T. Talkar \textit{et al.}, ``Assessment of speech and fine motor coordination in children with autism spectrum disorder,'' \textit{IEEE Access}, vol.~8, pp.~127535--127545, 2020.

\bibitem{ref7} M. Godel, F. Robain, F. Journal, N. Kojovic, K. Latr\`eche, G. Dehaene-Lambertz, and M. Schaer, ``Prosodic signatures of ASD severity and developmental delay in preschoolers,'' \textit{npj Digit. Med.}, vol.~6, Art. no.~99, 2023.

\bibitem{ref8} S. Sandoval, V. Berisha, R. L. Utianski, J. M. Liss, and A. Spanias, ``Automatic assessment of vowel space area,'' \textit{J. Acoust. Soc. Am.}, vol.~134, no.~5, pp.~EL477--EL483, 2013.

\bibitem{ref9} M. Abinath, S. M. Afaal, P. Anojan, R. Vithurshan, K. Pulasinga, and W. Tissera, ``Automatic speech recognition system to analyze autism spectrum disorder in young children,'' \textit{Int. J. Eng. Manag. Res.}, vol.~12, no.~5, pp.~1--8, 2022.

\bibitem{ref10} S. Bae \textit{et al.}, ``Multimodal AI for risk stratification in autism spectrum disorder: Integrating voice and screening tools,'' \textit{npj Digit. Med.}, vol.~8, Art. no.~01914-6, 2025.

\bibitem{ref11} Y. Zheng, S. Ma, X. Liu, L. Wang, M. Chen, J. Yuan, and H. Wang, ``Deep neural networks and deep deterministic policy gradient for early ASD diagnosis and personalized intervention in children,'' \textit{Sci. Rep.}, vol.~15, Art. no.~26854, 2025.

\bibitem{ref12} S. Saeedi, H. Bouraghi, M.-S. Seifpanahi, and M. Ghazisaeedi, ``Application of digital games for speech therapy in children: A systematic review of features and challenges,'' \textit{J. Healthcare Eng.}, vol.~2022, Art. no.~4814945, 2022.

\bibitem{ref13} L. Furlong, M. Morris, T. Serry, and S. Erickson, ``Mobile apps for treatment of speech disorders in children: An evidence-based analysis of quality and efficacy,'' \textit{PLoS ONE}, vol.~13, no.~8, Art. no.~e0201513, 2018.

\bibitem{ref14} A. E. Ortiz Castellanos, C.-M. Liu, and C. Shi, ``Deep mobile linguistic therapy for patients with ASD,'' \textit{Int. J. Environ. Res. Public Health}, vol.~19, no.~19, Art. no.~12857, 2022.

\bibitem{ref15} C.-H. Chueh \textit{et al.}, ``Implementation of a generative AI-powered digital interactive platform for clinical language therapy in children with language delay: A pilot study,'' \textit{Life}, vol.~15, no.~10, Art. no.~1628, 2025.

\bibitem{ref16} W. Katz \textit{et al.}, ``Opti-Speech: A real-time, 3D visual feedback system for speech training,'' in \textit{Proc. INTERSPEECH}, Singapore, 2014, pp.~1174--1178.

\bibitem{ref17} A. Aftab, C. A. Sehgal, M. M. Noohu, and G. Jaleel, ``Clinical effectiveness of AAC intervention in minimally verbal children with ASD: A systematic review,'' \textit{NeuroRegulation}, vol.~10, no.~4, pp.~239--252, 2023.

\bibitem{ref18} N. M. Nguyen, T. T. Nguyen, P.-N. Tran, C. P. Lim, N. T. Pham, and D. N. M. Dang, ``Multimodal fusion in speech emotion recognition: A comprehensive review of methods and technologies,'' \textit{Eng. Appl. Artif. Intell.}, vol.~163, Art. no.~112624, 2026.

\bibitem{ref19} A. Landowska, A. Karpus, T. Zawadzka, B. Robins, D. E. Barkana, H. Kose, T. Zorcec, and N. Cummins, ``Automatic emotion recognition in children with autism: A systematic literature review,'' \textit{Sensors}, vol.~22, no.~4, Art. no.~1649, 2022.

\bibitem{ref20} M. S. Remya, R. Raman, R. Sankaran, V. Namboodiri, and P. Nedungedi, ``Artificial intelligence for classification and enhancement of speech and language disorders: Techniques, applications, and future directions,'' \textit{IEEE Access}, vol.~13, pp.~1--28, 2025.

\bibitem{ref21} P. Govalkar, J. Fischer, F. Zalkow, and C. Dittmar, ``A comparison of recent neural vocoders for speech signal reconstruction,'' in \textit{Proc. 10th ISCA Speech Synthesis Workshop (SSW)}, Vienna, Austria, 2019, pp.~7--12.

\bibitem{ref22} J. Zhu, C. Yang, F. Samir, and J. Islam, ``The taste of IPA: Towards open-vocabulary keyword spotting and forced alignment in any language,'' in \textit{Proc. Annu. Meeting Assoc. Comput. Linguistics (ACL)}, 2024, pp.~1--16.

\bibitem{ref23} A. Barovic and A. Moin, ``TinyML for speech recognition,'' in \textit{Proc. IEEE 49th Annu. Comput., Softw., Appl. Conf. (COMPSAC)}, 2025, pp.~1--6.

\bibitem{ref24} A. Woubie and T. B\"{a}ckstr\"{o}m, ``Federated learning for privacy-preserving speaker recognition,'' \textit{IEEE Access}, vol.~9, pp.~149477--149487, 2021.

\end{thebibliography}

\end{document}